% !TeX root = ../main.tex

\chapter{Conclusiones y trabajo futuro}
\label{chap:conclusiones}

Este capítulo cierra la memoria recapitulando qué se ha obtenido y qué queda pendiente. La primera sección responde a los cuatro objetivos planteados en la introducción y valora en qué medida los resultados sostienen las dos hipótesis de trabajo, sin volver sobre el detalle de las tablas ni sobre la interpretación ya desarrollada en la discusión. La segunda recoge las líneas de continuación que se derivan del trabajo realizado, ordenadas según lo que cada una permitiría resolver de las preguntas que este estudio deja abiertas.

\section{Conclusiones}
\label{sec:conclusiones}

Este trabajo ha estudiado si atributos asociados a una empresa modifican las decisiones de inversión de comités multiagente basados en modelos de lenguaje. La conclusión principal es que el efecto observado no se manifiesta como un desplazamiento sistemático de las asignaciones a favor o en contra de los atributos sensibles. Donde aparece, lo hace como un aumento de la dispersión entre ejecuciones. Dos sistemas pueden presentar un \textit{gap} medio similar y, sin embargo, comportarse de forma muy distinta. Uno puede responder de manera consistente entre ejecuciones, mientras que otro puede mostrar una sensibilidad mucho más variable ante el atributo. Esta diferencia queda oculta cuando la evaluación se limita al promedio.

\textbf{O1.} El banco de pruebas construido mantiene constantes los fundamentales financieros y la posición de la empresa sujeto, sustituye los nombres reales por alias y dispone de dos referencias de variabilidad. El agente ciego proporciona el \textit{null} interno, al recibir entradas idénticas en ambas variantes de la comparación, y las parejas de cestas gemelas de la condición placebo constituyen el \textit{null} externo. La auditoría estructural verificó la coincidencia de las entradas ciegas en los 1.512 grupos examinados de las campañas de detección y mitigación. El placebo mostró además que aparecen diferencias entre ejecuciones incluso sin manipulación alguna, lo que obliga a interpretar cualquier efecto frente a un nivel de ruido y no frente al cero teórico.

\textbf{O2.} Los atributos manipulados no producen un desplazamiento medio generalizable, ya que ninguna de las 36 pruebas sobre el \textit{gap} medio conserva significación tras la corrección por comparaciones múltiples. La evidencia aparece en la dispersión: doce celdas interpretables presentan ratios de varianzas entre 3,10 y 8,22 respecto a la referencia interna. El fenómeno no se asocia de forma estable a determinadas cestas, puesto que ninguna de las 42 estimaciones de correlación intraclase sobrevive a la corrección, y tampoco aumenta al enriquecer las instrucciones. Donde sí aparece una separación clara es en la composición del comité: el exceso aparece en tres de las cuatro composiciones homogéneas y en ninguna de las dos heterogéneas. El contraste resulta además más limpio en génesis, antes de cualquier comunicación, lo que coincide con la localización temporal descrita por Nguyen \textit{et al.}~\cite{nguyen2025social}.

\textbf{O3.} La deliberación hace que las decisiones de los agentes dejen de ser independientes entre sí. En las seis composiciones, cuando los agentes visibles presentan una diferencia elevada entre control y variante, el agente ciego tiende a mostrar también una diferencia elevada en el turno siguiente. Este patrón aparece igualmente en la condición placebo, por lo que no puede atribuirse por sí solo al atributo manipulado. Solo \textit{homo\_mistral} presenta una diferencia significativa entre tratamiento y placebo tras la corrección por comparaciones múltiples ($q=0{,}026$), y esto se debe a una pendiente placebo raramente baja. Existe por tanto evidencia de una asociación diferencial en una composición, pero no de un mecanismo de propagación generalizable al resto de los comités.

\textbf{O4.} Las dos formas de mitigación examinadas ofrecen evidencias de distinto alcance. La composición heterogénea coincide con la ausencia total de exceso de dispersión en las condiciones evaluadas. Este resultado sugiere un posible efecto estabilizador de combinar modelos distintos, aunque no permite demostrar que la diversidad sea por sí sola la causa, puesto que en un comité homogéneo cada modelo ocupa los tres roles mientras que en uno mixto ocupa solo uno, y ambos factores varían a la vez. Debe considerarse por ello un candidato a mitigación estructural antes que una eficacia demostrada. La intervención por instrucciones, evaluada sobre \textit{homo\_qwen}, reduce la varianza visible en ocho de doce comparaciones y de forma consistente en geografía, si bien solo una de ellas excluye el valor de referencia con su intervalo de confianza. Ningún cambio en el \textit{gap} medio mantiene la significación tras la corrección, por lo que no hay evidencia de que las instrucciones hayan generado una sobrecorrección en la dirección contraria. La evidencia es prometedora, pero no concluyente, y las diferencias entre los diseños impiden determinar cuál de las dos estrategias resulta más eficaz.

Respecto a las hipótesis, los resultados apoyan parcialmente \textbf{H1}. El cambio aislado del atributo no provoca un desplazamiento medio sistemático, pero sí altera la dispersión en parte de las configuraciones, con la excepción de que esa segunda forma el efecto se identificó de manera exploratoria en \textit{homo\_llama} antes de replicarse en otras composiciones. La deliberación modifica el comportamiento observado, como sostenía \textbf{H2}, aunque los ratios posteriores a génesis no miden ya lo mismo, dado que el agente ciego deja de ser independiente de la información que producen sus compañeros.

La ausencia de una diferencia media no basta, en conjunto, para caracterizar la equidad de un sistema multiagente. Un atributo puede no favorecer ni perjudicar de forma sostenida a una empresa y volver, sin embargo, mucho menos estable el trato que recibe. Hacer visible esa forma de sensibilidad exige repeticiones, referencias de ruido y medidas de dispersión. Y su aparición depende de cómo esté compuesto el comité, lo que sitúa la arquitectura del sistema a la altura de los modelos que lo integran a la hora de explicar su comportamiento.

\section{Trabajo futuro}
\label{sec:trabajo_futuro}

La primera línea de continuación consiste en ampliar la cobertura experimental. Incorporar modelos más actuales, versiones de mayor tamaño y un número mayor de repeticiones permitiría comprobar hasta qué punto el exceso de dispersión depende de las arquitecturas concretas empleadas. El aumento de semillas resultaría especialmente útil para caracterizar efectos de menor magnitud y para determinar con mayor potencia si algunas cestas presentan una sensibilidad persistente, cuestión que el análisis actual no resuelve.

Los resultados de las composiciones heterogéneas abren la segunda línea. La ausencia de exceso en las configuraciones evaluadas apunta a que combinar modelos distintos estabiliza la decisión colectiva, pero el diseño no permite separar ese efecto del reparto concreto de modelos entre roles. Evaluar permutaciones de los mismos modelos entre puestos determinaría si la diversidad constituye por sí misma el mecanismo responsable. Un diseño en el que composiciones homogéneas y heterogéneas recibieran además las mismas instrucciones de mitigación permitiría comparar ambas vías bajo condiciones equivalentes y comprobar si resultan complementarias, algo que la evidencia actual no autoriza a decidir.

Una tercera línea afecta al eje de niveles de instrucción. El nivel más específico reproduce literalmente la descripción funcional del nivel anterior y le añade una única oración de identidad, de modo que la manipulación consiste exclusivamente en dotar al agente de una breve biografía. Esas biografías, sin embargo, sitúan a los tres agentes en centros financieros occidentales, y una de ellas declara además una especialización en narrativas de mercado occidentales. La comparación entre ambos niveles no distingue por ello el efecto de asignar una identidad del de asignar esa procedencia concreta, distinción especialmente relevante para el eje geográfico. Repetir el nivel con identidades equivalentes en extensión y detalle pero situadas en otras regiones permitiría separar ambos factores.

La dinámica deliberativa requiere igualmente una evaluación más amplia. Extender el número de rondas mostraría si el aumento de variabilidad posterior a la génesis continúa, se estabiliza o converge. Esa extensión permitiría además evaluar la formulación activa de mitigación, cuyo mecanismo de detección y corrección de influencias ajenas no puede operar en la decisión inicial. Esta es una de las principales preguntas que quedan abiertas, ya que la campaña de mitigación no se aplicó durante la deliberación y, por tanto, no permite saber si una intervención puede reducir la influencia que las decisiones de los agentes visibles ejercen sobre el agente ciego. Una campaña de mitigación durante los turnos de deliberación permitiría comprobar si la instrucción activa reduce la pendiente de transmisión sin reducir simultáneamente la del placebo, que es el contraste que da sentido a esa medida.

El banco de pruebas puede extenderse, por último, a otros atributos, contextos geográficos y tareas. Aplicar el mismo procedimiento a la concesión de crédito, apoyo a decisiones médicas, la selección de candidatos o el reparto de otros recursos mostraría si el exceso de dispersión es propio del escenario de inversión o una forma más general de sensibilidad en sistemas multiagente. Dentro del propio dominio financiero, incorporar resultados económicos posteriores permitiría estudiar si la inestabilidad observada ante cambios en el atributo se traduce en diferencias materiales de riesgo o rendimiento, cuestión que este trabajo deja expresamente fuera.

